\section{Global first-energy Green calculus and relative Aubin--Mabuchi geometry}
\label{sec:r11-global-first-energy}

Throughout this section the domain and background satisfy precisely
(TC1)--(TC2): $V$ is the strict original analytic Laurent box compactly
contained in the interior of one original affinoid domain, and the
arbitrary original symmetric middle-positive separately closed
$(p,p)$-current $S$ is defined on a neighborhood of its closed box.
Put $q=n-p\geq1$, retain the negative bounded barrier $\rho$ of
Lemma~\ref{lem:r11-strict-Laurent-barrier}, and abbreviate
$\mathcal E^1=\mathcal E^1_{S,\partial}(V)$.
The background is not assumed strongly positive. The positive energy
matrices used below come from rank-one original $(1,1)$ squares
wedged with the same middle-positive background.

The compact bounded Green identity is the frozen R4 input
\texttt{eq:relative-green-energy} in
\texttt{relative\_energy\_agent.tex}, with the analytic local
scope made explicit in Section~\ref{sec:r11-analytic-Laurent-collars}.
The bounded common-collar Aubin--Mabuchi variation is already frozen
in R3 \texttt{energy\_convex\_blocks.tex},
\texttt{lem:r3-energy-midpoint}. The new assertions here remove that
fixed compact collar by the proved unbounded zero-boundary energy
class. The complete argument keeps track of the global boundary
term; it does not assume that it vanishes merely because a potential
has pointwise zero boundary values.

\begin{lemma}[Global bounded Green exchange from finite curvature mass]
\label{lem:r11-global-bounded-Green}
Let $Q$ be any original symmetric middle-positive separately closed
$(n-1,n-1)$-current on $V$. Let $a,b\leq0$ be bounded classical
psh functions satisfying $\mathrm{BG}(a,0)$ and
$\mathrm{BG}(b,0)$. Suppose the two positive measures
$D(aQ)$ and $D(bQ)$ have finite total mass on $V$. Then
\[
 \begin{split}
  \int_V(-a)D(bQ)&=\int_V(-b)D(aQ),\\
  \int_V\Gamma_Q(a,b)&=\int_V(-a)D(bQ),\qquad
  \int_V\Gamma_Q(a)=\int_V(-a)D(aQ)<\infty.
 \end{split}
 \tag{GE1}
\]
The cross energy has finite total variation. These are identities
of the actual original Borel multipliers and energy measures.
\end{lemma}

\noindent\textbf{Status: \texttt{NEW\_PROVED}} as a global
zero-boundary extension of compact bounded Green exchange.

\begin{proof}
For $\epsilon>0$ put $a^\epsilon=\max(a,-\epsilon)$.
The bounded difference $a-a^\epsilon$ is supported in the compact
set $K_\epsilon=\overline{\{a<-\epsilon\}}\Subset V$.
Compact Green exchange therefore gives
\[
 \int(a-a^\epsilon)D(bQ)
     =\int bD(aQ)-\int bD(a^\epsilon Q).
 \tag{GE2}
\]
The measure $D(a^\epsilon Q)$ has the same finite total mass as
$D(aQ)$. Indeed their difference is $D((a^\epsilon-a)Q)$,
with compactly supported multiplier. A compact original smooth
cutoff equal to one on a neighborhood of that support has zero
pairing with this derivative. Outside the compact support the
two measures agree, by ordinary locality. This proves both
finiteness and equality of their total masses.

Fix $\delta>0$. The set
$L_\delta=\overline{\{b<-\delta\}}$ is compact in $V$.
Relative CLN for the fixed $Q$ and the normalized classical
function $(a^\epsilon+\epsilon)/\epsilon\in[0,1]$ gives
$D(a^\epsilon Q)(L_\delta)\leq C_{\delta,Q}\epsilon$.
Consequently
\[
 \left|\int bD(a^\epsilon Q)\right|
 \leq \|b\|_\infty C_{\delta,Q}\epsilon
                     +\delta D(aQ)(V).
 \tag{GE3}
\]
First let $\epsilon\downarrow0$, then $\delta\downarrow0$.
The last term in (GE2) tends to zero. Bounded convergence on
its left proves the first equality in (GE1).

The original bounded clipping rule, including the zero energy
of an exact scalar level, gives
$\Gamma_Q(a-a^\epsilon,a)
 =\mathbf1_{\{a<-\epsilon\}}\Gamma_Q(a)$.
Apply compact Green to this difference and $a$:
\[
 \int_{\{a<-\epsilon\}}\Gamma_Q(a)
                 =\int(a^\epsilon-a)D(aQ).
 \tag{GE4}
\]
As $\epsilon\downarrow0$ the right side increases to the
finite integral $\int(-a)D(aQ)$. The energy gives no mass
to $\{a=0\}$, so monotone convergence gives the diagonal
energy identity. Apply it also to $b$ and $a+b$; the latter
has the same boundary property and finite curvature mass.
Polarization and the first equality of (GE1) identify the
cross integral. The measure Cauchy--Schwarz inequality bounds
its total variation by the product of the square roots of the
two finite diagonal masses.
\end{proof}

For $u_0,\ldots,u_q\in\mathcal E^1$ define
\[
 \mathcal I_S(u_0,\ldots,u_q)
       =\int_V(-u_0)M(u_1,\ldots,u_q).
 \tag{GE5}
\]
It is finite by the already proved mixed estimate (PE11).
Its interpretation uses the actual potentials; all mixed
measures annihilate their joint pole union.

\begin{theorem}[Full symmetry of the global first mixed energy]
\label{thm:r11-global-first-mixed-symmetry}
The nonnegative functional (GE5) is symmetric in all $q+1$
slots and separately additive and positively homogeneous on
$\mathcal E^1$. In particular the undifferentiated slot may
be exchanged with any differentiated slot without a boundary
term. Every integral is over the whole ordinary domain $V$.
\end{theorem}

\noindent\textbf{Status: \texttt{NEW\_PROVED}} for the
unbounded first-power boundary class and arbitrary original $S$.

\begin{proof}
Use the decreasing finite-mass approximants from
Theorem~\ref{thm:r11-finite-mass-power-approximation}:
$b_{i,j}=\max(u_i,j\rho)$. For a fixed $j$ they all lie
in $[-j,0]$ and satisfy the compact-sublevel boundary condition.
Let $A_j=\sum_{i=0}^q b_{i,j}$. Since
$A_j\geq(q+1)j\rho$ and $\rho<0$, comparison with
$c(q+1)j\rho$ for $c>1$ has strict gap all of $V$ and gives
\[
          \mu_{A_j}(V)
              \leq[(q+1)j]^q\mu_\rho(V)<\infty
 \tag{GE6}
\]
after $c\downarrow1$. The positive mixed expansion of
$\mu_{A_j}$ shows that every mixed top measure formed from
these bounded slots has finite total mass.

Fix the last $q-1$ slots and use their bounded mixed product
as $Q$ in Lemma~\ref{lem:r11-global-bounded-Green}. Its two
curvature measures are finite by (GE6), so the lemma exchanges
the first two slots. Symmetry among the differentiated slots
is already part of bounded BT; hence the bounded mixed energy
is fully symmetric.

The approximants in each slot have the fixed lower barrier
$u_i\in\mathcal E^1$. The global critical mixed-energy
continuity of Theorem~\ref{thm:r11-critical-energy-barrier-continuity}
therefore passes the bounded identity to (GE5), on the whole
domain and with no escaped boundary mass. This proves full
symmetry. Additivity and positive homogeneity in the first
slot are ordinary integration identities; in the other slots
they follow from the unbounded mixed-current calculus. The
same symmetry transfers all of them to every slot.
\end{proof}

\begin{theorem}[Globally finite unbounded energy matrices]
\label{thm:r11-global-first-energy-matrix}
Let $a,b,c_2,\ldots,c_q\in\mathcal E^1$ and put
$Q=N(c_2,\ldots,c_q)$, with $Q=S$ when $q=1$.
The graded energy currents at the endpoint $\sigma=1$ are
finite measures on all of $V$ and satisfy
\[
 \begin{split}
 \int_V\Gamma_Q(a,b)
       &=\mathcal I_S(a,b,c_2,\ldots,c_q),\\
 \left|\int_V\Gamma_Q(a,b)\right|^2
       &\leq\mathcal I_S(a,a,c_2,\ldots,c_q)
                   \mathcal I_S(b,b,c_2,\ldots,c_q).
 \end{split}
 \tag{GE7}
\]
Every cross energy has finite total variation, and the same
matrix positivity holds for finite real linear combinations
of the displayed potentials. Thus it applies to differences
even though a difference need not be psh.
\end{theorem}

\noindent\textbf{Status: \texttt{NEW\_PROVED}} as global
finiteness and Green identification; local endpoint existence
is the separate proved theorem (PG10)--(PG12).

\begin{proof}
Fix $\epsilon>0$ and put $a^\epsilon=\max(a,-\epsilon)$.
This bounded function belongs to $\mathcal E^1$ by solidity.
For integers $j>\epsilon$, write $a_j=\max(a,-j)$ and
let $Q_j$ use the same canonical level $j$ in its other slots.
The difference $a_j-a^\epsilon$ is supported in the one
compact set $K_\epsilon=\overline{\{a<-\epsilon\}}$.
The bounded compact Green formula and clipping locality give
\[
 \int_{\{a<-\epsilon\}}\Gamma_{Q_j}(a_j)
                =\int(a^\epsilon-a_j)D(a_jQ_j).
 \tag{GE8}
\]
Insert a fixed original compact cutoff equal to one near
$K_\epsilon$ when reading this identity. The endpoint graded
energy convergence (PG11) passes its left side in local total
variation. Its indicator is independent of $j>\epsilon$.
The weighted curvature convergence (PG4) with $a=1,k=q$
passes its right side, since $a+k=q+1=q+\sigma$ for
$\sigma=1$. Include the bounded slot $a^\epsilon$ in the
fixed family; its canonical truncation is itself for large $j$.
Thus
\[
 \int_{\{a<-\epsilon\}}\Gamma_Q(a)
           =\int(a^\epsilon-a)M(a,c_2,\ldots,c_q).
 \tag{GE9}
\]
The energy gives zero mass to the joint pole set by (PG12).
It also gives no mass to $\{a=0\}$: the corresponding bounded
energies vanish there, and (PG11) is local total-variation
convergence on the identical set $\{a_j=0\}=\{a=0\}$.
Let $\epsilon\downarrow0$ in (GE9). Both sides increase,
and the right side has finite limit by (PE11). This proves
the global diagonal identity in (GE7).

Apply it to $a+b$ as well. The full symmetry from the
preceding theorem and polarization give the cross identity.
The positive local energy matrix and an increasing compact
exhaustion give the global Cauchy--Schwarz inequality and the
total-variation bound. No convergence of global unweighted
energy masses was assumed in order to obtain their finiteness.
\end{proof}

\begin{theorem}[The constant-one first-energy Holder inequality]
\label{thm:r11-optimal-first-energy-Holder}
For all $u_0,\ldots,u_q\in\mathcal E^1$,
\[
 \boxed{\quad
 \mathcal I_S(u_0,\ldots,u_q)
                    \leq\prod_{i=0}^q e_1(u_i)^{1/(q+1)}.
 \quad}
 \tag{GE10}
\]
In particular
\[
 e_1(u+v)^{1/(q+1)}
           \leq e_1(u)^{1/(q+1)}+e_1(v)^{1/(q+1)}.
 \tag{GE11}
\]
The constant one in (GE10) is optimal whenever a potential
of positive first energy exists. This is an integral inequality;
it makes no pointwise Alexandrov--Fenchel or hyperbolicity
claim for the arbitrary middle-positive background.
\end{theorem}

\noindent\textbf{Status: \texttt{NEW\_PROVED}}. The earlier
(PE11) supplied a finite constant for every power and is used
to justify integrability before this sharper first-power proof.

\begin{proof}
The zero-energy case already follows from (PE11), including
its proved zero-factor limiting argument. Otherwise rescale
each potential so that $e_1(u_i)=1$. Put $d=q+1$.
Among the finitely many mixed values
\[
 B_\alpha=\mathcal I_S(
             \underbrace{u_0,\ldots,u_0}_{\alpha_0},\ldots,
             \underbrace{u_q,\ldots,u_q}_{\alpha_q}),
       \qquad \alpha_i\in\mathbf N,\quad\sum_i\alpha_i=d,
\]
choose a largest value $M$, and among its indices choose one
maximizing $\sum_i\alpha_i^2$. If this index has two positive
entries $\alpha_i,\alpha_l$, the two-slot energy
Cauchy--Schwarz inequality (GE7), with all remaining slots
fixed, gives
\[
       M^2=B_\alpha^2
           \leq B_{\alpha+e_i-e_l}B_{\alpha-e_i+e_l}
           \leq M^2.
\]
Since a diagonal value is one, $M\geq1>0$, and both
neighboring values must equal $M$. At least one neighboring
index has a strictly larger sum of squares: the two changes
are $2(\alpha_i-\alpha_l)+2$ and its opposite plus $2$.
This contradicts the choice. Thus a maximizing index is
pure, and $M=e_1(u_i)=1$. Every mixed value is at most one,
proving (GE10) after undoing the normalization.

Expand $e_1(u+v)$ by the fully symmetric multilinearity of
(GE5), and apply (GE10) term by term. The binomial theorem
gives (GE11). If all slots equal one potential with positive
energy, equality holds in (GE10), proving optimality.
\end{proof}

\begin{proposition}[The relative capacity quotient determines every original product]
\label{prop:r11-capacity-representative-quotient}
Let $u_i,v_i\in\mathcal F_{S,\partial}(V)$ for a finite
family of indices, and suppose
\[
       C_S(\{u_i\ne v_i\};V)=0\quad\hbox{for every }i.
 \tag{GE12}
\]
Then every mixed original product of the $u_i$ is equal to
the corresponding product of the $v_i$, in every degree
$0\leq k\leq q$, against every original complementary test
and after every actual moment refinement. For one pair,
\[
 \mu_u=\mu_v
  \quad\Longleftrightarrow\quad C_S(\{u\ne v\};V)=0
  \quad\Longleftrightarrow\quad C_S^*(\{u\ne v\};V)=0.
 \tag{GE13}
\]
If one member of a pair in (GE13) belongs to
$\mathcal E^\sigma_{S,\partial}$, so does the other, with
the same power energy. All defined graded energy currents
and integrable weighted products agree after these replacements.
\end{proposition}

\noindent\textbf{Status: \texttt{NEW\_PROVED}} as a
representative quotient within the specified actual-potential
classes. It is not the analytic PS/PL relation quotient GQ.

\begin{proof}
Let $E$ be the finite union of the exceptional Borel sets in
(GE12); it has zero raw capacity. At each canonical level $j$,
the bounded differences $u_{i,j}-v_{i,j}$ vanish outside $E$.
Every bounded mixed background $Q_j$ over $S$ gives $E$ zero
variation in each original coefficient by the global
coefficient bound (DG3), including degree zero. Therefore
$(u_{i,j}-v_{i,j})Q_j=0$ as an original current. The bounded
telescoping formula
\[
 \prod_{i=1}^k Du_{i,j}\wedge S-
 \prod_{i=1}^k Dv_{i,j}\wedge S
       =\sum_{i=1}^k D\bigl((u_{i,j}-v_{i,j})Q_{i,j}\bigr)
 \tag{GE14}
\]
is zero; $Q_{i,j}$ uses the already replaced slots before
$i$ and the unreplaced slots after it. Pass the uncut local
coefficient total-variation limits from the all-degree theorem.
This proves product equality without assuming a locality
principle on an arbitrary equality Borel set.

For a single pair, equality of the top measures gives both
one-sided outer-capacity domination statements from
Corollary~\ref{cor:r11-boundary-outer-domination}. Their
boundary-gap hypotheses hold because both potentials are
negative and have compact negative sublevels. Countable
subadditivity gives outer capacity zero of their inequality
set. Outer zero implies raw zero, and the proved product
equality gives the remaining implication in (GE13).

All products annihilate $E$ in their coefficient variations,
again by (DG3) and the increasing cut construction. Thus
the actual multiplier values of the two families agree
almost everywhere for each such coefficient. This preserves
integrability, power energies and weighted products.
The graded energy formula (PG10) then proves energy-current
equality wherever its degree hypotheses are satisfied.
\end{proof}

\begin{theorem}[Global Aubin--Mabuchi variation and energy rigidity]
\label{thm:r11-global-Aubin-Mabuchi}
For $u\in\mathcal E^1$ put
\[
             \mathscr A_S(u)=\frac1{q+1}\int_V u\,d\mu_u
                              =-\frac{e_1(u)}{q+1}.
 \tag{GE15}
\]
For $u,v\in\mathcal E^1$, let $h=v-u$ and
$u_t=(1-t)u+tv$, $0\leq t\leq1$. Then $\mathscr A_S(u_t)$
is a polynomial and
\[
 \begin{split}
 \frac d{dt}\mathscr A_S(u_t)&=\int_V h\,d\mu_{u_t},\\
 \frac {d^2}{dt^2}\mathscr A_S(u_t)
              &=-q\int_V\Gamma_{N(u_t^{q-1})}(h)\leq0.
 \end{split}
 \tag{GE16}
\]
All displayed signed integrals are absolutely convergent.
The functional is increasing under pointwise order, and
\[
 \int(v-u)\,d\mu_v
    \leq\mathscr A_S(v)-\mathscr A_S(u)
    \leq\int(v-u)\,d\mu_u.
 \tag{GE17}
\]
The symmetric discrepancy
\[
 \begin{split}
 \mathscr I_S(u,v)
    &:=\int_V(u-v)(d\mu_v-d\mu_u)\\
    &=\sum_{l=0}^{q-1}
       \int_V\Gamma_{N(u^l,v^{q-1-l})}(u-v)\geq0
 \end{split}
 \tag{GE18}
\]
vanishes exactly when the equivalent conditions (GE13) hold.
Consequently $\mathscr A_S$ is strictly concave on the
relative capacity quotient: equality with its chord at one
strict interior parameter holds exactly for equivalent
potentials. No pointwise separation or positivity on every
ordinary open set is assumed for the background.
\end{theorem}

\noindent\textbf{Status: \texttt{NEW\_PROVED}} as the
global unbounded energy extension and quotient rigidity.
The bounded common-collar variation formula itself is an
\texttt{INTERNAL\_INPUT}, as specified above.

\begin{proof}
Full symmetric multilinearity of (GE5) makes the energy a
polynomial along the segment. Differentiating its $q+1$
equal slots gives the first formula in (GE16). Differentiating
once more gives
$q\int hD(hN(u_t^{q-1}))$. The global energy matrix identity
(GE7), extended by bilinearity to the difference $h$, is
exactly its negative in (GE16). Finiteness follows from
(GE10) in each expanded mixed term. Concavity follows from
matrix positivity, and the tangent inequalities (GE17) follow
by integrating the decreasing first derivative. If $u\leq v$
the first derivative is nonnegative, proving monotonicity.

The telescoping identity for $\mu_v-\mu_u$ and the same
global Green identity give (GE18). Suppose its sum is zero.
Every summand is nonnegative, so for each
$Q_l=N(u^l,v^{q-1-l})$ one has
$\Gamma_{Q_l}(h)(V)=0$.
For an original compact smooth scalar test $\varphi$, choose
an original compact cutoff $\chi$ equal to one on its
differential support, and use canonical bounded approximants
$h_j=v_j-u_j$ and $Q_{l,j}$. Compact integration by parts
and energy Cauchy--Schwarz give
\[
 |\langle D(h_jQ_{l,j}),\varphi\rangle|^2
       \leq\left(\int\chi\Gamma_{Q_{l,j}}(h_j)\right)
             \left(\int\chi\Gamma_{Q_{l,j}}(\varphi)\right).
 \tag{GE19}
\]
The first factor tends to zero by the endpoint local
energy total-variation convergence (PG11). The second is
bounded: its scalar differentials are fixed original smooth
forms, and $Q_{l,j}$ has local coefficient total-variation
convergence. The weighted convergence (PG4) also gives
$h_jQ_{l,j}\to hQ_l$, so (GE19) proves $D(hQ_l)=0$
against every original test. Summing the recursive telescoping
formula yields $\mu_u=\mu_v$. Conversely this equality
makes the defining expression in (GE18) zero. Proposition
\ref{prop:r11-capacity-representative-quotient} gives precisely
the other equivalences.

If the energy agrees with its chord at one interior point,
concavity forces it to be affine on the whole segment.
Its endpoint derivatives are equal, so (GE16) gives
$\mathscr I_S(u,v)=\mathscr A_S'(0)-\mathscr A_S'(1)=0$.
Conversely equivalent potentials have identical mixed products
and energies by the preceding proposition, and the energy is
constant on their segment. This proves the claimed strict
concavity on the quotient.
\end{proof}

\begin{corollary}[The variational necessary condition retains its domain gate]
\label{cor:r11-variational-necessary-condition}
Let $\mu$ be a positive Radon datum such that
$\int(-v)\,d\mu<\infty$ for every $v\in\mathcal E^1$.
If $u\in\mathcal E^1$ solves $\mu_u=\mu$, then it maximizes
\[
                   v\longmapsto\mathscr A_S(v)-\int v\,d\mu
 \tag{GE20}
\]
on $\mathcal E^1$. Conversely, a maximizer satisfies the
current equation provided that for every original compact
smooth scalar $\varphi$ there is $\epsilon_\varphi>0$
such that $u+t\varphi\in\mathcal E^1$ for all
$|t|<\epsilon_\varphi$. This last variation hypothesis is
explicit and is not asserted for a general classical psh
potential or arbitrary datum.
\end{corollary}

\noindent\textbf{Status: \texttt{DIRECT\_COROLLARY}} for
necessity and \texttt{CONDITIONAL} for the stated converse.

\begin{proof}
The upper tangent inequality in (GE17) gives necessity.
Under the printed two-sided variation hypothesis, apply
the first derivative identity to the admissible segment
$u+t\varphi$. Maximality for both signs gives
$\int\varphi\,d\mu_u=\int\varphi\,d\mu$ for every
original compact smooth scalar test. This is exactly the
original current equation. No envelope differentiability or
unstated admissibility of arbitrary smooth perturbations is
used.
\end{proof}

\paragraph{Scope and rejected shortcut.}
The tempting inference from first-energy dual boundedness to
an unrestricted ambient classical Dirichlet solution remains
\texttt{FALSE}, by Corollary~\ref{cor:r11-dual-energy-insufficient}.
Concavity alone supplies no compactness or representability
theorem for a maximizing family. The two-sided variation gate
in the last corollary cannot be discarded by analogy with a
smooth strictly elliptic problem. The relative capacity quotient
in (GE12)--(GE14) identifies actual potentials in the proved
boundary classes; it does not identify analytic PS/PL
presentations, does not enlarge the class of delta tests, and
does not close GQ or low-real-degree CTO.

One possible alternate approximation is
$\max(u_i-\epsilon_j,j\rho)$, which imposes a common ordinary
collar. That route requires separate uniform control of the
small-value energy near the moving collar. The proof above
instead uses (GE3) and the already proved decreasing finite-mass
approximation; no claim about the alternate route is needed.
